\section{Bounded forcing on an expanding array}
\label{sec:bounded}\label{cbe:section}

Diffusion can also be controlled by enlarging the region that carries
each computational state.  We now use this option to keep every mixed
derivative of the force bounded.  A horizontal incompressible field
installs all possible transitions between finitely many addresses at
each stage.  A positive vertical impulse selects one address.  The
vertical velocity then obeys an advection--diffusion equation, and a
moving cutoff measures the mass that remains near the selected
computation.  The radii grow fast enough that the cumulative diffusion
loss is less than a fixed detection margin.

We prove the second part of Theorem~\ref{thm:main}, on
$\Omega=\R^2\times\T$ with the solution class
\eqref{eq:comparison-class}. The horizontal force support will be
compact on each finite interval, while the scalar has diffusion tails.
We verify those tails in the exact comparison spaces before applying
uniqueness.

\subsection{Finite addresses and their transitions}
\label{cbe:addresses}

The only computational input is the injective local compiler of
\cite[Lemma~3.1]{stationary2026}.  Its output is
a finite partial one-head table with states $i=1,\ldots,N$, tape
alphabet $\Sigma$, and an initial tape differing from the full blank
symbol at finitely many cells.  It has one terminal state, with no outgoing rule.
Its initialized run reaches that state exactly when the original
machine halts; otherwise every successive transition is defined.
Every target state fixes the displacement of its incoming rules, and
the target state together with the full written symbol identifies
the incoming rule.  The compiler first adds a nonterminal initial
step, so these assertions include an input on which the original
machine is already halted.  Its history mechanism is the reversible
recording method used by Bennett~\cite{Bennett1973}; here we use
the stated finite table, including its full-domain incoming-rule
property.

Let $b=|\Sigma|\ge2$ and assign its symbols distinct digits
$d_\alpha\in\{0,\ldots,b-1\}$, with the full blank symbol assigned
zero.  The integer
$\sum_{j\ge0}d_{\alpha_j}b^j$ encodes an eventually blank stack whose
top is $\alpha_0$.  Zero tails are interpreted as infinitely many
blanks.  For a configuration, let $x$ encode the tape strictly left
of the head, nearest cell first, and let $y$ encode the head cell and
the tape to its right.  If the rule reads $\alpha$, writes $\beta$
and enters state $i'$, then $d_\alpha=y\bmod b$, and the new stacks are
\begin{equation}\label{cbe:stack}
\begin{array}{c|cc}
\text{move}&x'&y'\\ \hline
\mathsf R&bx+d_\beta&\lfloor y/b\rfloor\\
\mathsf N&x&y-d_\alpha+d_\beta\\
\mathsf L&\lfloor x/b\rfloor&
 d_\gamma+bd_\beta+b^2\lfloor y/b\rfloor,
 \quad d_\gamma=x\bmod b.
\end{array}
\end{equation}
These formulas push and remove exactly the indicated tape symbols.

They also give an injective map on all configurations where a rule
is defined.  The target state determines which row applies.  In the
right row the lowest digit of $x'$ recovers $d_\beta$; in the stationary
row the lowest digit of $y'$ does so; in the left row the second
lowest digit of $y'$ does so, while its lowest digit recovers
$d_\gamma$.  The compiler then identifies the old state and read
symbol.  The inverses are, respectively,
\[
 (x,y)=\bigl((x'-d_\beta)/b,\,by'+d_\alpha\bigr),\qquad
 (x,y)=\bigl(x',\,y'-d_\beta+d_\alpha\bigr),
\]
and
\[
 (x,y)=\bigl(bx'+d_\gamma,\,
                  b\lfloor y'/b^2\rfloor+d_\alpha\bigr).
\]
This proves injectivity without restricting to the initialized run.

Choose $m\ge1$ effectively so that the two initial stacks are smaller
than $b^m$, and define
\begin{equation}\label{cbe:counts}
 B_n=b^{m+n},\qquad K_n=NB_n^2=K_0D^n,\qquad D=b^2.
\end{equation}
Every update with $0\le x,y<B_n$ has $0\le x',y'<B_{n+1}$.
For the largest expression, in the left row,
\[
 y'\le(b-1)+b(b-1)+b^2(B_n/b-1)=bB_n-1.
\]
All other bounds follow directly from \eqref{cbe:stack}.
At level $n$, give the triple $(i,x,y)$ the address
\begin{equation}\label{cbe:address}
 k=1+x+B_ny+B_n^2(i-1),\qquad 1\le k\le K_n.
\end{equation}
Integer division recovers the triple.  Thus every address with a
defined instruction determines one target address $k'$ at level
$n+1$, and distinct defined addresses have distinct targets.
Only this one-step lookup is used to prescribe the force.

\subsection{Disjoint translation gates at increasing scales}
\label{cbe:gates}

We now assign each address a square and move all defined source
squares to their targets simultaneously.  Three successive motions
keep the squares separated: descent along source columns, horizontal
motion along private rows, and vertical motion along target columns.

Use the smooth step $\theta$ from Section~\ref{sec:compact-blocks}, and put
\[
 \Theta(r)=\theta(2r-1/2).
\]
The function $\Theta$ makes its transition between
$1/4$ and $3/4$.  Introduce the scales
\begin{equation}\label{cbe:scales}
\begin{aligned}
 C_*&=3000,&\Lambda&=4D,
 &R_0&=1024C_*(1+\nu)\Lambda(K_0D+2),\\
 R_n&=R_0\Lambda^n,&S_n&=16R_n,
 &T_n&=S_{n+1}(K_{n+1}+2),\\
 t_0&=1,&t_{n+1}&=t_n+T_n.
\end{aligned}
\end{equation}
The constant $C_*$ will bound the Laplacian of a cutoff at radius
$R$ by $C_*/R^2$. For scalar mass at most one, a stage of duration
$T_n$ will therefore lose at most $\nu C_*T_n/R_n^2$ from that
cutoff. The radii grow faster than the number of addresses so that
these losses form a summable series. The durations $T_n$ are long
enough for the most distant gate to move at bounded speed.
In particular $R_n,T_n\ge1$, and $t_n\ge n+1$.

Let $s_i=1$ for the terminal state and $s_i=-1$ for every other state.
An address $k$ of state $i$ has center
\[
 p_{n,k}=(S_nk,s_iS_n).
\]
The known initial configuration determines a center $p_*$ at level
zero.  It lies on the negative row.  For a defined source address $k$
at level $n$, with target $k'$ in state $i'$, put
\[
 Y_k=-(k+2)S_n,\qquad
 \theta_j(t)=\Theta\bigl(3(t-t_n)/T_n-(j-1)\bigr),
 \quad j=1,2,3.
\]
Its center path on $[t_n,t_{n+1}]$ is
\begin{equation}\label{cbe:path}
\begin{aligned}
 c_{n,k,1}(t)&=(1-\theta_2(t))S_nk
                         +\theta_2(t)S_{n+1}k',\\
 c_{n,k,2}(t)&=-(1-\theta_1(t))S_n
            +(\theta_1(t)-\theta_3(t))Y_k
                         +\theta_3(t)s_{i'}S_{n+1}.
\end{aligned}
\end{equation}
It starts at $p_{n,k}$ because every defined source is nonterminal,
and it ends at $p_{n+1,k'}$.  The supports of
$\dot\theta_1,\dot\theta_2,\dot\theta_3$ occur in successive
disjoint time slots, with stationary collars between them.

During the first slot, distinct centers have first coordinates
separated by at least $S_n$.  During the second, their private
heights $Y_k$ are separated by at least $S_n$.  During the third,
their target first coordinates are separated by at least
$S_{n+1}$, by the injectivity just proved.  The same separations
hold through the intervening collars.  If the target is nonterminal,
every vertical segment stays at or below the source row:
\begin{equation}\label{cbe:negative}
 c_{n,k,2}(t)\le-S_n\qquad(t_n\le t\le t_{n+1}).
\end{equation}

To realize one center path, set
\[
 \chi(a)=\prod_{\ell=1}^2\theta(a_\ell+3)\theta(3-a_\ell),
 \qquad \nabla^\perp=(\partial_{X_2},-\partial_{X_1}).
\]
For $c=c_{n,k}$, $G=\dot c$ and $\xi=X-c(t)$, define
\begin{equation}\label{cbe:gate}
 W_{n,k}(t,X)=\nabla^\perp
   \left[\chi(\xi/R_n)(G_1\xi_2-G_2\xi_1)\right].
\end{equation}
This field is divergence free, equals $G$ when
$|\xi|_\infty\le2R_n$, and vanishes when
$|\xi|_\infty\ge3R_n$.  The gate supports are disjoint at each
time, since their centers are separated in one coordinate by at
least $16R_n>6R_n$.
Figure~\ref{fig:translation-gates} displays the separation coordinate
in each time slot.

\begin{figure}[ht]
\centering
\begin{tikzpicture}[x=1cm,y=.8cm,>=Latex,font=\scriptsize,
  route/.style={line width=.7pt,->},
  first/.style={blue!65!black},
  second/.style={green!45!black},
  gateone/.style={draw=blue!65!black,fill=blue!12},
  gatetwo/.style={draw=green!45!black,fill=green!12}]
\begin{scope}
  \node[font=\small] at (2,1.35) {1. Source columns};
  \draw[dashed,gray] (0,0)--(4.2,0);
  \node[anchor=south west] at (0,0) {$X_2=0$};
  \draw[route,first] (.7,-.55)--(.7,-1.65);
  \draw[route,second] (3.5,-.55)--(3.5,-2.5);
  \fill[first] (.7,-.55) circle (1.5pt);
  \fill[second] (3.5,-.55) circle (1.5pt);
  \node[anchor=east] at (3.3,-.55) {source row};
  \draw[gateone] (.56,-1.24) rectangle (.84,-.96);
  \draw[gatetwo] (3.36,-1.665) rectangle (3.64,-1.385);
  \draw[<->] (.7,-2.85)--(3.5,-2.85);
  \node[below] at (2.1,-2.85) {$\ge S_n$};
\end{scope}
\begin{scope}[xshift=5.1cm]
  \node[font=\small] at (2,1.35) {2. Private rows};
  \draw[dashed,gray] (0,0)--(4.2,0);
  \draw[route,first] (.7,-1.65)--(2.8,-1.65);
  \draw[route,second] (3.5,-2.5)--(1.3,-2.5);
  \fill[first] (.7,-1.65) circle (1.5pt);
  \fill[second] (3.5,-2.5) circle (1.5pt);
  \draw[gateone] (1.61,-1.79) rectangle (1.89,-1.51);
  \draw[gatetwo] (2.26,-2.64) rectangle (2.54,-2.36);
  \draw[<->] (3.95,-1.65)--(3.95,-2.5);
  \node[anchor=west] at (3.95,-2.075) {$\ge S_n$};
\end{scope}
\begin{scope}[xshift=10.2cm]
  \node[font=\small] at (2,1.35) {3. Target columns};
  \draw[dashed,gray] (0,0)--(4.2,0);
  \draw[route,first] (2.8,-1.65)--(2.8,.8);
  \draw[route,second] (1.3,-2.5)--(1.3,-.8);
  \fill[first] (2.8,.8) circle (1.5pt);
  \fill[second] (1.3,-.8) circle (1.5pt);
  \node[anchor=west] at (2.95,.8) {terminal};
  \node[anchor=east] at (1.15,-.8) {nonterminal};
  \draw[gateone] (2.66,-.565) rectangle (2.94,-.285);
  \draw[gatetwo] (1.16,-1.79) rectangle (1.44,-1.51);
  \draw[<->] (1.3,-2.85)--(2.8,-2.85);
  \node[below] at (2.05,-2.85) {$\ge S_{n+1}$};
\end{scope}
\end{tikzpicture}
\caption{Two installed gates in the three successive motions of one
stage (schematic, not to scale). Small squares bound instantaneous
supports. Source columns, private rows and target columns separate
simultaneous supports, although routes from different phases may cross.
Nonterminal routes stay below $X_2=0$; terminal routes ascend across it.
The squares do not represent the support of the diffusing scalar.}
\label{fig:translation-gates}
\end{figure}

Let $W=0$ on $[0,1]$, and on $[t_n,t_{n+1}]$ sum
\eqref{cbe:gate} over every level-$n$ address with a defined rule.
The sum is finite.  Every gate is zero near the time-slot endpoints,
so these formulas join smoothly, including at $t=1$.  The resulting
field has compact horizontal support on each finite time interval
and has the exact plateau property
\begin{equation}\label{cbe:plateau}
 W(t,X)=\dot c_{n,k}(t)
 \quad\hbox{whenever }|X-c_{n,k}(t)|_\infty\le R_n.
\end{equation}

Here is why the growing array costs no derivative bound.  Every
coordinate used in \eqref{cbe:path} is bounded by a fixed multiple
of $T_n$: for example,
$|Y_k|\le S_n(K_n+2)\le T_n$ and
$S_{n+1}k'\le T_n$.  Therefore, for every $j\ge1$,
\begin{equation}\label{cbe:center-bounds}
 |c_{n,k}^{(j)}(t)|\le C_j T_n^{1-j},
\end{equation}
with $C_j$ independent of $n$ and $k$.  Differentiating the potential
in \eqref{cbe:gate} writes the gate as
$A((X-c)/R_n)\dot c$ for a fixed smooth compactly supported matrix
$A$.  A spatial derivative contributes $R_n^{-1}$; time derivatives
contribute derivatives of $c$ divided by $R_n$, together with higher
derivatives of $c$.  Since $R_n,T_n\ge1$, the product and chain rules
and \eqref{cbe:center-bounds} bound each fixed mixed derivative
uniformly in $n,k$.  At most one gate contributes at a point.
The same global bounds thus hold for $W$.

\subsection{A prescribed impulse and its scalar evolution}
\label{cbe:scalar-section}

The horizontal field is now fixed.  The remaining force is a unit
vertical impulse near $p_*$.  Prescribe
\begin{equation}\label{cbe:force}
\begin{aligned}
 f_h&=\partial_tW+(W\cdot\nabla_X)W-\nu\Delta_XW,\\
 g(t,X)&=\theta'(t)\prod_{\ell=1}^2\theta'\bigl(X_\ell-(p_*)_\ell+1/2\bigr),\\
 f(t,X,z)&=(f_h(t,X),g(t,X)).
\end{aligned}
\end{equation}
All terms are prescribed independently of the unknown vertical
velocity.  The bounds for $W$ prove global boundedness of all mixed
derivatives of $f$, and its horizontal support is compact on each
finite time interval.  Moreover, $g\ge0$, its support lies in
$0\le t\le1$ and $|X-p_*|_\infty\le1/2$, and
\begin{equation}\label{cbe:impulse}
 \int_{\mathbb R^2}g(t,X)\,dX=\theta'(t).
\end{equation}

Define $\rho$ by solving
\begin{equation}\label{cbe:scalar}
 \partial_t\rho+W\cdot\nabla_X\rho
       =\nu\Delta_X\rho+g,\qquad \rho(0,X)=0.
\end{equation}
We need integrability as well as smoothness, because both the mass
test and the noncompact uniqueness argument use it.

\begin{lemma}[Smooth scalar evolution with integrable tails]
\label{cbe:scalar-lemma}
Equation \eqref{cbe:scalar} has a global smooth solution.  Every
space-time derivative is continuous in time with values in
$C_0(\mathbb R^2)\cap L^1(\mathbb R^2)$, where $C_0$ denotes
continuous functions tending to zero at infinity.  Furthermore,
\begin{equation}\label{cbe:mass}
 \rho\ge0,\qquad
 \mathcal M(t):=\int_{\mathbb R^2}\rho(t,X)\,dX=\theta(t)\le1.
\end{equation}
\end{lemma}

\begin{proof}
For $s>0$, let
$H_s(X)=(4\pi\nu s)^{-1}\exp(-|X|^2/(4\nu s))$.
Since $\operatorname{div}_XW=0$, the mild equation is
\begin{equation}\label{cbe:mild}
 \rho(t)=\int_0^t H_{t-r}*g(r)\,dr
 -\sum_{\ell=1}^2\int_0^t
       (\partial_\ell H_{t-r})*(W_\ell(r)\rho(r))\,dr.
\end{equation}
For an integer $a\ge0$, let $E_a$ be the Banach space of $C^a$
functions whose derivatives of order at most $a$ lie in $C_0\cap L^1$,
with norm
\[
 \|v\|_{E_a}=\sum_{|\alpha|\le a}
       \bigl(\|\partial^\alpha v\|_\infty
                    +\|\partial^\alpha v\|_1\bigr).
\]
Heat convolution is a contraction on each derivative norm and is
strongly continuous on $E_a$.  The latter follows from continuity
of translations in $C_0$ and $L^1$ and the approximate-identity
property of the heat kernels.  Multiplication by $W_\ell(t)$ is a
bounded operator on $E_a$, continuous in time in operator norm on
finite intervals, by the bounded derivatives of $W$.

The elementary estimate
$\|\partial_\ell H_s\|_1\le C_\nu s^{-1/2}$ shows that the linear
operator in the second term of \eqref{cbe:mild} has norm at most
$C_a\sqrt{\tau}$ on $C([0,\tau];E_a)$.  For sufficiently small
$\tau$ this is a contraction.  The same integrable kernel bound
proves continuity of the time integrals at their upper endpoints.
On a later interval one adds heat evolution of the value at its
left endpoint.  A uniform short interval length on each $[0,T]$
therefore constructs a solution on every finite interval.  Uniqueness
in $E_0$ identifies the solutions constructed for the different $a$.

Put $F=g-\operatorname{div}_X(W\rho)$.  The spatial regularity just
proved gives $F\in C([0,T];E_a)$ for every $a$, and
\eqref{cbe:mild} becomes
$\rho(t)=\int_0^tH_{t-r}*F(r)\,dr$.
For $v\in E_{a+2}$,
\[
 \frac{d}{ds}(H_s*v)=\nu H_s*\Delta v
\]
extends continuously to $s=0$ in $E_a$.  Differentiation under the
integral gives $\partial_t\rho=F+\nu\Delta\rho$ in every $E_a$.
Repeated differentiation of this identity proves the claimed time
regularity.  In particular all smoothness statements hold one-sided
at $t=0$; the source is flat there.

All terms of \eqref{cbe:mild} are absolutely integrable in space and
time.  Integrating it and using
$\int\partial_\ell H_s=0$ gives
$\mathcal M(t)=\int_0^t\theta'(r)\,dr=\theta(t)$.
For positivity, fix $T$ and set
$v_\varepsilon=\rho+\varepsilon(1+t)$.
Continuity into $C_0$ implies uniform decay of $\rho(t,X)$ as
$|X|\to\infty$ for $0\le t\le T$: the image of this compact time
interval is compact in $C_0$ and can be covered by finitely many
small sup-norm balls.  Hence $v_\varepsilon$ is positive outside
one fixed compact set.  It is also positive initially.  At a first
zero, attained inside that compact set, its time derivative is
nonpositive, its gradient vanishes and its Laplacian is nonnegative.
But its equation has source $g+\varepsilon>0$, a contradiction.
Letting $\varepsilon\downarrow0$ proves $\rho\ge0$.
\end{proof}

We can now verify the fluid equation and its solution class.  Set
\begin{equation}\label{cbe:solution}
 u(t,X,z)=(W(t,X),\rho(t,X)),\qquad p(t,X,z)=0.
\end{equation}
The divergence vanishes and $u\cdot\nabla=W\cdot\nabla_X$ on
these $z$-independent components.  Thus the horizontal equation is
the definition of $f_h$ in \eqref{cbe:force}, and the vertical
equation is \eqref{cbe:scalar}.  Both components start from zero.

On $[0,T]$, the horizontal field and its derivatives have support in
one compact set.  Every derivative of $\rho$ is continuous in both
$L^1$ and $L^\infty$ by Lemma~\ref{cbe:scalar-lemma}.
The inequality
\[
 \|v\|_2^2\le\|v\|_1\|v\|_\infty
\]
applied also to differences at two times gives continuity into
$L^2$.  Applying it to the spatial derivatives through order two
and to $\partial_t\rho$ proves exactly the velocity regularity in
\eqref{eq:comparison-class}; the same lemma gives the required bounded velocity
and gradient.  Pressure zero belongs to the stated pressure class.

The comparison argument of Section~\ref{sec:scalar-structure} now
applies in exactly the established class. In particular, its pressure
cutoff error uses $L^2$ pressure, and its velocity errors use the scalar's
integrable tails, not compact support of the scalar. It gives uniqueness
among all three-dimensional competitors in
\eqref{eq:comparison-class}. The half-plane integral is well defined by
Lemma~\ref{cbe:scalar-lemma}.

\subsection{A moving cutoff and the detection margin}
\label{cbe:detection}

We have constructed the force and its unique solution without
following the selected computation.  Only to prove the halting
equivalence do we now follow that run.  During loading put $c(t)=p_*$.
Thereafter concatenate the paths \eqref{cbe:path} corresponding to
its successive configurations, up to the first terminal endpoint
if one occurs.  Each required gate is present because of the
integer-stack bounds and the compiler's continuation property.
The paths agree at endpoints and are stationary in endpoint collars.

An explicit test function will quantify the mass retained near
$c(t)$.  On $[0,1]$ put
$P(s)=10s^3-15s^4+6s^5$, extending it by zero to the left and one
to the right.  This extension is $C^2$, takes values in $[0,1]$,
and satisfies $|P''|\le360$.  Define
\[
 A_*(s)=P(2(1+s))P(2(1-s)),\qquad
 \varphi_R(a)=A_*(a_1/R)A_*(a_2/R).
\]
Then $0\le\varphi_R\le1$, it equals one on
$|a|_\infty\le R/2$, and it vanishes for $|a|_\infty\ge R$.
The transition intervals of the two factors in $A_*$ are disjoint;
on each, the other factor equals one.  Consequently
$|A_*''|\le1440$, and
\begin{equation}\label{cbe:laplacian}
 \|\Delta\varphi_R\|_\infty\le2880R^{-2}\le C_*R^{-2}.
\end{equation}

Take $R=R_0$ during loading and $R=R_n$ during the $n$th
transition, and within each such interval put
\[
 \mathcal I(t)=\int_{\mathbb R^2}
            \varphi_R(X-c(t))\rho(t,X)\,dX.
\]
Differentiating under the integral and integrating by parts against
the compactly supported $C^2$ cutoff gives
\begin{equation}\label{cbe:moving}
\begin{aligned}
 \mathcal I'(t)
  ={}&\int_{\mathbb R^2}
       \bigl((W-\dot c)\cdot\nabla\varphi_R(X-c)
                    +\nu\Delta\varphi_R(X-c)\bigr)\rho\,dX\\
   &+\int_{\mathbb R^2}\varphi_R(X-c)g\,dX.
\end{aligned}
\end{equation}
The transport term vanishes exactly by \eqref{cbe:plateau}; during
loading both $W$ and $\dot c$ are zero.  Positivity,
$\mathcal M\le1$, and \eqref{cbe:laplacian} bound the diffusion
term below by $-\nu C_*R^{-2}$.  The loading source is supported
where $\varphi_{R_0}=1$, because $R_0\ge1$; it contributes
$\mathcal M'(t)$.  After loading the source vanishes.

There is no loss at a radius change.  Indeed
$R_{n+1}\ge2R_n$, so the new cutoff equals one on the entire support
of the previous one, with the same center at the joint.  Nonnegativity
then implies that $\mathcal I$ cannot decrease at that change.
Starting with zero mass and integrating \eqref{cbe:moving} yields,
at every time through the first terminal endpoint,
\begin{equation}\label{cbe:loss}
 \mathcal I(t)\ge\mathcal M(t)-\delta,\qquad
 \delta=\nu C_*\left(R_0^{-2}
                   +\sum_{n=0}^{\infty}\frac{T_n}{R_n^2}\right)
       <\frac1{24}.
\end{equation}
The last inequality follows directly from the chosen scales:
\[
 \frac{T_n}{R_n^2}=
 \frac{16\Lambda}{R_0}
      \left(K_0D(D/\Lambda)^n+2\Lambda^{-n}\right).
\]
Here $D/\Lambda=1/4$ and $\Lambda\ge4$, so each geometric series
has sum at most $4/3$.  Therefore
\[
 \nu C_*\sum_{n\ge0}\frac{T_n}{R_n^2}
 \le\frac{\nu}{1+\nu}\frac{64}{3\cdot1024}<\frac1{48}.
\]
Also $R_0\ge1024C_*(1+\nu)$ gives
\[
 \frac{\nu C_*}{R_0^2}
 \le\frac{\nu}{1024^2C_*(1+\nu)^2}<\frac1{48}.
\]
This proves the strict bound in \eqref{cbe:loss}.

At every completed level before halting, more than $3/4$ of the
unit mass lies in the address square of radius $R_n$ about the
correct center.  At a terminal endpoint the preceding cutoff radius
already gives that conclusion, and its square lies in the larger
level-$n$ square.  Other level-$n$ address squares are disjoint
from this one, so none can contain mass $1/2$.  Thus the velocity
field records the intermediate configurations as well as their
eventual outcome.

If the run reaches a terminal center at level $n$, its height is
$S_n=16R_n$.  The cutoff at arrival lies wholly in $X_2>0$ and
captures at least $1-\delta>3/4$ of the mass.  Its arrival time is
finite, so \eqref{eq:plane-event} holds.  This also handles an initially
halted original input: the compiler starts nonterminal and reaches
its terminal state after the finite preliminary simulation.
No assertion about the captured mass after the first terminal
endpoint is needed for this existential event.

If the machine never halts, the loading center is on the negative
row and every transition has a nonterminal target.  By
\eqref{cbe:negative}, the moving cutoff remains entirely in
$X_2<0$ at every physical time.  Hence
\[
 \int_{\{X_2>0\}}\rho(t,X)\,dX
 \le\mathcal M(t)-\mathcal I(t)
 \le\delta<\frac1{24}<\frac14.
\]
The vertical period has length one, so this integral equals the
one in \eqref{eq:plane-event}.  The threshold $1/2$ therefore separates
the two cases with a strict margin.

\subsection{Effective finite prescriptions}
\label{cbe:effective}

Every integer, scale, address lookup and single-rule update above
is computed directly from the finite machine table and input.
For any finite interval $[0,T]$, choose an integer $Q>T$.
Since $t_n\ge n+1$, only the levels $0\le n<Q$ can contribute
there.  Extend each gate by zero through its time collars.  Evaluation
on $[0,T]$ then uses a finite sum over those levels and their finite
address sets; it requires no decision about which side of a
computable real switching time contains the argument.

The smooth profiles and every requested derivative are effectively
evaluable. Near a flat endpoint, derivatives of $e^{-1/r}$ are
$e^{-1/r}$ times explicit polynomials in $1/r$.  The estimate
$s^d e^{-s}\le(d+k)!s^{-k}$ supplies an effective error bound there,
and the denominator in $\theta$ is bounded below by $e^{-2}$.
Finite differentiation of the gate formulas gives effective
bounds of every specified derivative order.  Formula
\eqref{cbe:force} is consequently a finite program for the force
and its derivatives to any prescribed accuracy.

The force installs every defined one-step transition. It never follows
the distinguished computation to decide which gate to install. The
initial impulse selects its configuration; the subsequent solution
carries the mass through the prescribed array. This completes the
proof of the second part of Theorem~\ref{thm:main}.

\subsection{Effective evaluation of the scalar}
\label{cbe:scalar-effective}

The scalar can also be evaluated effectively on each finite time
interval, including its derivatives and integrable tails. This is a
stronger computability property, not an input to the force prescription.
We give the error control because pointwise computation alone would
not justify restarting the scalar solver on this unbounded domain.
Fix a derivative
order $a$ and a time bound $T$. Let $M_a$ be an effective bound for
multiplication by each $W_\ell(t)$ on $E_a$, and let $C_\nu$ be a
constant large enough to include the sum of both derivative-kernel
bounds. On a step of length $h$ choose
\[
 q=2C_\nu M_a\sqrt h\le\tfrac12.
\]
Then the Volterra operator $B$ in the second term of
\eqref{cbe:mild}, with the lower time limit replaced by the step's
initial time, has norm at most $q$. Write $y$ for the sum of heat
evolution of the initial datum and the source integral on that step.
The truncation $\sum_{j=0}^N B^j y$ has error at most
\[
 \frac{q^{N+1}}{1-q}\,\|y\|_{C_tE_a}.
\]
An initial-data error $\epsilon_0$ and an additional error
$\epsilon_1$ in the integral equation, measured in $C_tE_a$, change
its solution by at most $(\epsilon_0+\epsilon_1)/(1-q)$.
All the norms used here have computable upper bounds from the finite
coefficient formulas and the preceding step. Thus choosing $N$ and
the quadrature accuracies gives a prescribed step error. Finitely
many steps cover $[0,T]$, so their error budgets can be chosen
backwards to achieve any specified final accuracy.

For completeness, the approximants carry a computable spatial-tail
bound for every derivative through order $a$, in both supremum and
$L^1$ norm. The source $g$ and each product $W_\ell v$ have known
compact support on $[0,T]$, even when the approximant $v$ does not.
In a derivative-kernel integral, split the time lag at $\eta>0$.
If $\|v\|_{C_tE_a}\le A$, the part with lag less than $\eta$ has
$E_a$ norm at most $2C_\nu M_a A\sqrt\eta$.
For lags in $[\eta,T]$, derivatives of the explicit Gaussian have
computable supremum and $L^1$ tails, uniform in that interval.
Convolution with data of known compact support and known $E_a$ norm
therefore has both tail bounds outside the support radius plus a
computably chosen radius. The source integral is treated in the
same way, with its bounded heat-kernel operator in place of the
integrable derivative-kernel singularity.

At a restart, the exact previous scalar is not assumed compactly
supported. Carry forward its certified $E_a$ approximant and error.
Multiply that approximant by a smooth cutoff equal to one on a large
ball. Its certified derivative tails bound the cutoff error in $E_a$
by the product rule. Heat convolution contracts that error, while
the compactly supported truncation has explicit Gaussian tails,
uniformly from lag zero to $h$. This supplies the homogeneous term
at the next step with both its $E_a$ error and its tail bounds.
Starting from zero, induction proves that every finite Picard
approximation and every endpoint datum has the required effective
representation. On the remaining compact sets, the computably smooth
integrands can be integrated with derivative error bounds. A uniform
error for each derivative on the retained spatial ball bounds its
supremum error there and its $L^1$ error by the ball's area times that
error; the certified tails control the complement. Time derivatives
are recovered from \eqref{cbe:scalar}.
